Este trabalho investiga a predição de morte súbita cardíaca em pacientes com
insuficiência cardíaca crônica a partir do eletrocardiograma (ECG) Holter de
24~horas. Os dados provêm da base MUSIC (PhysioNet) e
compreendem sinais de três derivações ortogonais (X, Y e Z) amostrados a
200~Hz, além de variáveis clínicas de 992 pacientes de hospitais espanhóis. Adotou-se um desenho caso-controle restrito a
pacientes em ritmo sinusal, em que os casos correspondem a óbitos por morte
súbita cardíaca em até quatro anos e os controles, a sobreviventes com
seguimento completo de quatro anos, totalizando 517 indivíduos e 55 eventos.
Quatro abordagens foram comparadas sob o mesmo protocolo de validação cruzada
estratificada em cinco partições, com divisão no nível do paciente, limiar de
decisão definido na validação interna e intervalos de confiança por
reamostragem \textit{bootstrap}. A primeira aplica redes neurais profundas
diretamente ao sinal --- convolucionais, residuais, com atenção e de fusão ---, em dezenove configurações com 14~mil a
16,5~milhões de parâmetros. A segunda extrai marcadores do sinal ---
variabilidade da frequência cardíaca, morfologia do batimento e ectopia ---
em segmentos de 5~minutos e de 30~segundos. A
terceira utiliza apenas variáveis clínicas. Na quarta, uma rede aprende a
prever atributos do ecocardiograma, do ECG e do laboratório, e sua
representação é usada na predição. As redes neurais obtiveram área
sob a curva ROC (AUC) entre 0,48 e 0,56, compatíveis com o acaso. Os marcadores alcançaram AUC entre
0,59 e 0,62, e as variáveis clínicas, 0,66 (IC~95\,\%: 0,58--0,74). A
combinação das variáveis clínicas com os marcadores, com seleção feita dentro
de cada partição, chegou a 0,68 (IC~95\,\%: 0,60--0,76), valor próximo ao
relatado na literatura para a mesma base, mas o ganho sobre as variáveis
clínicas isoladas não foi estatisticamente significativo. A representação da
rede orientada por atributos clínicos alcançou, sozinha, 0,68 (IC~95\,\%:
0,60--0,76), equivalente às variáveis clínicas, mas não acrescentou informação
aos exames medidos, cujo melhor modelo, com o ecocardiograma, chegou a 0,73.
Os resultados
indicam que, com o número de eventos disponível, não foi possível aprender
do sinal bruto um padrão preditivo de morte súbita treinando a rede
diretamente para o desfecho, que a informação prognóstica do sinal coincide
com a dos exames já disponíveis, e que estimativas mais
altas dependem de escolhas de avaliação --- seleção de variáveis sobre a
coorte inteira, divisões únicas dos dados, tratamento de segmentos como
observações independentes e inclusão de ritmos não sinusais ---, cujos
efeitos este trabalho quantifica sobre os mesmos pacientes.

\textbf{Palavras-chave:} Morte Súbita Cardíaca. ECG Holter. Aprendizado
Profundo. Variabilidade da Frequência Cardíaca. Insuficiência Cardíaca.
Estratificação de Risco.
